\documentclass[11pt]{article}
\usepackage[a4paper,margin=1in]{geometry}
\usepackage{amsmath,amssymb,bm,booktabs}

\title{Equinoctial Short-Period Map: Alternative Constructions\\
\large Notes on removing the near-circular/equatorial degradation}
\author{Internal note}
\date{March 2026}

\begin{document}
\maketitle

\section{Problem Statement}

The current GEqOE short-period map constructs corrections in the polar
variables $(g, \Psi, Q, \Omega)$ extracted from the equinoctial pairs
$(p_1, p_2)$ and $(q_1, q_2)$. While the equinoctial pairs themselves are
regular at $g \to 0$ (circular) and $Q \to 0$ (equatorial), the polar
extraction $\Psi = \arctan(p_1/p_2)$ and $\Omega = \arctan(q_1/q_2)$
introduces angular variables that are poorly determined in these limits.
The short-period corrections $\delta\Psi$ and $\delta\Omega$ develop large
amplitudes that, when transformed back to equinoctial components
via $\delta p_1 = g\cos\Psi\,\delta g + \ldots$, produce amplified
along-track position errors.

This note outlines four approaches to eliminate this limitation, followed by
a feasibility study (Section~\ref{sec:feasibility}) that confirms
Approach~1 is viable and identifies the concrete implementation path.

\section{Approach 1: Direct Equinoctial Short-Period Map}

Construct short-period corrections directly in $(p_1, p_2, q_1, q_2)$
space rather than $(g, \Psi, Q, \Omega)$.

The $F$-variable residue decomposition framework is element-agnostic:
the forcing functions $\dot p_1, \dot p_2, \dot q_1, \dot q_2$ in the
GEqOE ODE (Eqs.~45--51 of Ba\"u et al.) are already expressed in
equinoctial form. The homological equation
\[
\frac{\partial u_{1}^{p_j}}{\partial M} = \dot p_j - \langle \dot p_j \rangle_M
\]
is solved by exactly the same Laurent-polynomial / residue technique,
since $\dot p_j$ has the same rational structure in $F$ (same denominator
$(F+q)^2(1+qF)^2$, different numerator polynomial). This eliminates the
polar-form singularity entirely.

\textbf{Advantages:}
\begin{itemize}
\item No $\arctan$ extraction or reconstruction needed.
\item Corrections remain bounded as $g \to 0$ or $Q \to 0$.
\item The residue method scales identically (same pole structure).
\end{itemize}

\textbf{Cost:} Re-derive and re-validate all short-period kernels. The
symbolic pipeline already handles arbitrary forcing polynomials, so the
implementation effort is modest.

\section{Approach 2: Historical Precedent}

Brouwer's original theory (1959) was singular at $e = 0$ and $i = 0$;
reformulation in non-singular Delaunay variables by Lyddane (1963) and
in equinoctial variables by Aksnes (1972) removed these singularities.
The GEqOE averaged theory is ripe for a similar evolution: the averaging
framework (near-identity transformation, residue decomposition) is
element-agnostic---only the choice of short-period correction variables
matters.

The key insight is that the \emph{mean-element ODE} is already well-behaved
(the coefficient functions depend on $g$, $Q$, not on $\Psi$ or $\Omega$
individually). The singularity is confined entirely to the short-period
reconstruction layer.

\section{Approach 3: Hybrid Approach}

Use $(g, \omega)$ for the secular rates---where they are well-conditioned
and yield compact expressions---but $(p_1, p_2)$ directly for the
short-period corrections.

The secular rates involve only $|m| \le n-2$ harmonics of $\omega$,
which are smooth functions of $(p_1, p_2)$:
\[
\cos(m\omega) = T_m(\cos\omega) - \ldots, \qquad
\omega = \arctan\frac{p_1}{p_2} - \arctan\frac{q_1}{q_2}.
\]
The Chebyshev expansion in $\cos\omega$ and $\sin\omega$ remains bounded
and can be evaluated without extracting $\omega$ explicitly.

\textbf{Advantage:} Reuses existing mean-element expressions while
improving only the short-period layer.

\section{Approach 4: Regularized Angle Extraction (Band-Aid)}

Use $\operatorname{atan2}$-based extraction with explicit quadrant handling
and singular-limit smoothing:
\[
\delta p_1 = \delta g \cdot \frac{p_1}{g} + g \cdot \delta\Psi \cdot \frac{p_2}{g},
\qquad
\delta p_2 = \delta g \cdot \frac{p_2}{g} - g \cdot \delta\Psi \cdot \frac{p_1}{g}.
\]
In the limit $g \to 0$, both $\delta g$ and $g \cdot \delta\Psi$ remain
$O(\varepsilon)$, so the product $g \cdot \delta\Psi$ is well-behaved even
though $\delta\Psi$ itself diverges. The key is to evaluate the products
$g\cos\Psi \cdot \delta\Psi$ and $g\sin\Psi \cdot \delta\Psi$ directly
from the symbolic expressions, never computing $\delta\Psi$ in isolation.

\textbf{Advantage:} Minimal code change---restructure the evaluation order.

\textbf{Disadvantage:} Does not address the fundamental issue; merely avoids
the worst numerical artifacts.

%% ======================================================================
\section{Feasibility Study: Complex-Variable Formulation}
\label{sec:feasibility}
%% ======================================================================

Approach~1 has been validated numerically
(\texttt{scripts/equinoctial\_sp\_feasibility.py}).
The key construction uses the complex eccentricity and complex inclination
to reduce the equinoctial SP problem to the existing residue machinery with
no new infrastructure.

\subsection{Complex eccentricity and inclination}

Define the complex eccentricity and inclination vectors
\[
\zeta \;=\; p_2 + i\,p_1 \;=\; g\,e^{i\Psi},
\qquad
\eta \;=\; q_2 + i\,q_1 \;=\; Q\,e^{i\Omega}.
\]
Their time derivatives factorize as
\[
\dot\zeta \;=\; (\dot g + i\,g\,\dot\Psi)\,e^{i\Psi}
         \;=\; (\dot g + i\,g\,\dot\Psi)\,w\,e^{i\Omega},
\qquad
\dot\eta \;=\; (\dot Q + i\,Q\,\dot\Omega)\,e^{i\Omega},
\]
where $w = e^{i\omega}$ and $\Psi = \omega + \Omega$.

\subsection{Reduced rate series}

The ``reduced'' complex rates (with the frozen factor $e^{i\Omega}$ stripped)
\begin{align}
\frac{\dot\zeta_{\mathrm{red}}}{\nu\lambda_n}
  &= \left[\frac{\dot g}{\nu\lambda_n}
     + i\,g\,\frac{\dot\Psi}{\nu\lambda_n}\right] w, \label{eq:zeta_red}\\[4pt]
\frac{\dot\eta_{\mathrm{red}}}{\nu\lambda_n}
  &= \frac{\dot Q}{\nu\lambda_n}
     + i\,Q\,\frac{\dot\Omega}{\nu\lambda_n}, \label{eq:eta_red}
\end{align}
are finite Laurent polynomials in $(w, F)$ with coefficients that are
rational functions of $(q, Q)$ alone.

\textbf{Critical cancellation.}
The $\dot\Psi$ dimensionless rate contains a term
$-(n-1)/(2g\beta)\,D^{-n}P_n(\hat z)(F + F^{-1})$, singular as $g \to 0$.
When multiplied by $g$ in~\eqref{eq:zeta_red}, this becomes the bounded
expression $-(n-1)/(2\beta)\,D^{-n}P_n(\hat z)(F + F^{-1})$, since
$g = 2q/(1+q^2)$ and $1/g = (1+q^2)/(2q)$ cancel algebraically.
The resulting Laurent coefficients are regular rational functions
of~$q$ at $q = 0$.

\subsection{Pipeline integration}

The reduced series~\eqref{eq:zeta_red}--\eqref{eq:eta_red} are constructed
by algebraic combination of the existing polar rate series:
\[
\texttt{zeta\_series} = \texttt{\_poly\_mul}\!\Big(
  \texttt{\_poly\_add}\big(\texttt{g\_series},\;
  \texttt{\_poly\_scale}(\texttt{Psi\_series},\, i\,g)\big),\;
  \{(1,0): 1\}
\Big).
\]
The multiplication by $\{(1,0): 1\}$ shifts the $w$-harmonic index by $+1$.
This Laurent polynomial is then processed by the \emph{same}
\texttt{integrate\_harmonic\_residue} solver used for the polar variables:
no new residue kernels, no new integration routines.

Reconstruction at evaluation time applies the frozen rotation:
\[
\delta p_1 = \operatorname{Im}\!\big(\delta\zeta_{\mathrm{red}}
  \cdot e^{i\Omega}\big),
\qquad
\delta p_2 = \operatorname{Re}\!\big(\delta\zeta_{\mathrm{red}}
  \cdot e^{i\Omega}\big),
\]
and analogously for $(\delta q_1, \delta q_2)$ from $\delta\eta_{\mathrm{red}}$.

\subsection{Numerical validation}

\paragraph{Algebraic equivalence.}
For moderate eccentricity ($g = 0.05$) and inclination ($Q = 0.36$), the
equinoctial and polar SP corrections agree to machine precision
($|\Delta(\delta p_j)| \sim 10^{-19}$) for both $J_2$ and $J_3$. The two
routes are algebraically identical.

\paragraph{Regularity at $g \to 0$.}
All $\zeta$-kernel Laurent coefficients converge to finite values as
$q \to 0$, confirmed by symbolic limit evaluation. The $J_2$ kernel
has 3~harmonics ($m \in \{-1, 1, 3\}$) and the $J_3$ kernel has
4~harmonics ($m \in \{-2, 0, 2, 4\}$), all finite at the circular limit.

\paragraph{Stability sweep.}
Evaluating the $J_2$ short-period correction $|\delta p_1|$ across seven
decades of eccentricity:

\medskip
\begin{center}
\begin{tabular}{ccc}
\toprule
$g$ & $|\delta\Psi|$ (polar) & $|\delta p_1|$ (equinoctial) \\
\midrule
$10^{-1}$  & $1.3 \times 10^{-4}$  & $1.32 \times 10^{-4}$ \\
$10^{-3}$  & $1.9 \times 10^{-3}$  & $1.30 \times 10^{-4}$ \\
$10^{-5}$  & $1.7 \times 10^{-1}$  & $1.30 \times 10^{-4}$ \\
$10^{-7}$  & $1.7 \times 10^{+1}$  & $1.30 \times 10^{-4}$ \\
\bottomrule
\end{tabular}
\end{center}
\medskip

\noindent
The polar $\delta\Psi$ grows as $O(1/g)$ while the equinoctial $\delta p_1$
remains constant at $O(\varepsilon)$, confirming that the singularity is
eliminated by construction.

%% ======================================================================
\section{Plan Forward}
\label{sec:plan}
%% ======================================================================

Based on the feasibility results, Approach~1 via the complex-variable
formulation is the selected path. The existing polar SP map is retained
as-is; the equinoctial route is built in parallel.

\subsection*{Stage 1: Full kernel generation (J2--J5)}

Extend the feasibility script to generate and cache equinoctial SP
coefficients for all four zonal degrees. Store in
\texttt{generated\_coefficients.py} alongside the existing polar data under
new keys (\texttt{zeta}, \texttt{eta}).

\subsection*{Stage 2: Parallel evaluator}

Add \texttt{mean\_to\_osculating\_state\_equinoctial} (scalar and batch) to
\texttt{short\_period.py}. This function:
\begin{enumerate}
\item Evaluates $\delta\zeta_{\mathrm{red}}$ and $\delta\eta_{\mathrm{red}}$
      from the cached Laurent kernels.
\item Applies the $e^{i\Omega}$ rotation to recover
      $(\delta p_1, \delta p_2, \delta q_1, \delta q_2)$.
\item Adds corrections to the mean state directly---no polar extraction,
      no $\arctan$, no angle reconstruction.
\item Recovers $K_{\mathrm{osc}}$ from the generalized Kepler equation
      (unchanged from the polar route).
\end{enumerate}
The $\delta M$ correction is shared with the polar route (it has no
polar-coordinate dependence).

\subsection*{Stage 3: Validation}

Run the full \texttt{extended\_validation.py} suite with the equinoctial
evaluator and compare against the polar route and Cowell ground truth.
The critical test cases are:
\begin{itemize}
\item Near-circular LEO ($e = 0.001$): expect along-track error reduction
      from $\sim\!1.5$~km to $\sim\!0.5$~km (matching the moderate-$e$ baseline).
\item Near-equatorial ($i = 5^\circ$): expect similar improvement in the
      inclination-plane error.
\item Retrograde and high-eccentricity cases: expect no regression.
\end{itemize}

\subsection*{Stage 4: Paper integration}

If validation confirms the improvement:
\begin{itemize}
\item Remove the near-circular/equatorial caveat from Section~A.2.
\item Add a brief note on the complex-variable construction (the formulation
      is compact enough for an appendix paragraph).
\item Update Table~12 with the improved error metrics.
\end{itemize}

\section{Recommendation (Updated)}

Approach~1 via the complex-variable formulation
(Section~\ref{sec:feasibility}) is confirmed feasible and is the selected
path forward. The construction reuses the entire existing residue pipeline,
requires no new symbolic infrastructure, and eliminates the $g \to 0$ and
$Q \to 0$ singularities by algebraic cancellation rather than numerical
workaround.

\subsection*{Log correction in the equinoctial route}

The log correction $C_{\log}\,\varphi(K)$ from the polar route carries
over identically to the equinoctial formulation.  The function
$\varphi = \arctan\!\bigl(q\sin f\,/\,(1 + q\cos f)\bigr)$ depends only
on the fast angle $f$ and the eccentricity parameter~$q$, so it is shared
between the polar and equinoctial evaluation paths.

For the $\zeta$-reduced rate (eccentricity channel), the log coefficient
$C_{\log}^{\zeta}$ is constructed from the same Taylor coefficient $E_1$
as the polar route but applied to the combined
$\dot g + ig\dot\Psi$ integrand.  Crucially, the equinoctial log
coefficients inherit the $g \to 0$ regularity of the equinoctial
formulation: $g \cdot E_1^{\Psi}$ remains bounded as $g \to 0$ (the $1/g$
pole in $\dot\Psi$ is cancelled by the $g$ prefactor in
$\zeta_{\mathrm{red}} = g e^{i\Psi}$), so $C_{\log}^{\zeta}$ is regular
at $g = 0$, just as the rational part is.

\end{document}
